\documentclass[12pt]{article}

\usepackage{sbc-template}
\usepackage{graphicx,url}
\usepackage[utf8]{inputenc}
\usepackage[brazil]{babel}
\usepackage{booktabs}
\usepackage{multirow}
\usepackage{array}
\usepackage{amsmath}
\usepackage{pgfplots}
\pgfplotsset{compat=1.18}

\sloppy

\title{Benchmarking de CPU, Mem\'oria e Armazenamento em Sistemas\\
com Arquitetura de Mem\'oria Unificada: Uma Avalia\c{c}\~ao Experimental}

\author{Luiza Barbosa Almeida da Silva\inst{1}}

\address{Centro de Estudos e Sistemas Avan\c{c}ados do Recife (CESAR School)\\
  Recife -- PE -- Brasil
  \email{lbas2@cesar.school}
}

\begin{document}

\maketitle

\begin{abstract}
Modern computing platforms based on ARM architectures with Unified Memory
Architecture (UMA) represent a significant paradigm shift in hardware design,
integrating CPU, GPU, and memory into a single System-on-Chip (SoC). Despite
growing adoption, empirical benchmarking studies on these platforms in academic
contexts remain scarce. This paper presents a controlled experimental evaluation
of a UMA-based system, measuring CPU, storage, and memory performance using
Cinebench R23, Blackmagic Disk Speed Test 3.0, and sysbench 1.0.20. Three
repetitions were performed per benchmark, with mean, standard deviation, and
95\% confidence interval reported. Results show single-core CPU performance of
$2156.67 \pm 6.81$ pts (CI$_{95\%}$: [2139.76, 2173.58]), positioning the
system above the Intel Core i7-1165G7 (1,532 pts) listed in the Cinebench R23
public ranking, although broader comparisons with current desktop x86
processors require further data. Sequential storage read reached $2809.1$
MB/s (median), and unified memory
bandwidth achieved $32{,}285$ MB/s in read with $0.03$ ms average latency.
Findings demonstrate that UMA platforms deliver competitive performance with
high single-core efficiency, while thermal throttling under sustained
multi-threaded load introduces measurable variability in multi-core scores.
All raw data, scripts, and reproduction instructions are publicly available.
\end{abstract}

\begin{resumo}
  Plataformas computacionais modernas baseadas em arquiteturas ARM com
  \textit{Unified Memory Architecture} (UMA) representam uma ruptura
  significativa no paradigma de design de hardware, integrando CPU, GPU e
  mem\'oria em um \'unico \textit{System-on-Chip} (SoC). Apesar da crescente
  ado\c{c}\~ao, estudos emp\'iricos de benchmarking dessas plataformas em
  contexto acad\^emico s\~ao escassos. Este artigo apresenta uma avalia\c{c}\~ao
  experimental controlada de um sistema baseado em UMA, medindo o desempenho
  de CPU, armazenamento e mem\'oria por meio das ferramentas Cinebench R23,
  Blackmagic Disk Speed Test 3.0 e sysbench 1.0.20. Foram realizadas tr\^es
  repeti\c{c}\~oes por benchmark, reportando-se m\'edia, desvio-padr\~ao e
  intervalo de confian\c{c}a de 95\%. Os resultados indicam desempenho
  single-core de $2156{,}67 \pm 6{,}81$ pts (IC$_{95\%}$: [2139,76;
  2173,58]), posicionando o sistema acima do Intel Core i7-1165G7 (1.532 pts)
  listado no ranking p\'ublico do Cinebench R23, embora compara\c{c}\~oes mais
  amplas com processadores x86 desktop atuais demandem dados adicionais.
  A leitura sequencial de
  armazenamento atingiu $2809{,}1$ MB/s (mediana), e a largura de banda de
  mem\'oria unificada alcan\c{c}ou $32.285$ MB/s em leitura com lat\^encia
  m\'edia de $0{,}03$ ms. Os achados demonstram que plataformas UMA entregam
  desempenho competitivo com alta efici\^encia single-core, enquanto o
  \textit{throttling} t\'ermico sob carga multi-thread sustentada introduz
  variabilidade mensur\'avel nas pontua\c{c}\~oes multi-core. Todos os dados
  brutos, scripts e instru\c{c}\~oes de reprodu\c{c}\~ao est\~ao
  publicamente dispon\'iveis.
\end{resumo}

\noindent\textbf{Palavras-chave:} benchmarking; arquitetura de mem\'oria
unificada; desempenho de hardware; Apple Silicon; avalia\c{c}\~ao experimental.


%% ============================================================
\section{Introdu\c{c}\~ao}
%% ============================================================

A avalia\c{c}\~ao objetiva de desempenho de hardware constitui uma
compet\^encia central na forma\c{c}\~ao de profissionais de tecnologia da
informa\c{c}\~ao. Especifica\c{c}\~oes t\'ecnicas nominais fornecidas por
fabricantes frequentemente n\~ao refletem com fidelidade o comportamento dos
sistemas sob cargas de trabalho reais, tornando o \textit{benchmarking} ---
medi\c{c}\~ao comparativa de desempenho por meio de testes padronizados e
reprodut\'iveis --- a principal ferramenta metodol\'ogica dispon\'ivel para
avalia\c{c}\~ao emp\'irica de hardware \cite{lilja2005measuring}.

A transi\c{c}\~ao de parte da ind\'ustria de semicondutores para arquiteturas
ARM de alto desempenho, com a integra\c{c}\~ao de CPU, GPU e mem\'oria em um
\'unico SoC sob o paradigma de \textit{Unified Memory Architecture} (UMA),
representa uma mudan\c{c}a estrutural no design de sistemas computacionais
\cite{hennessy2017computer}. Nessa arquitetura, CPU e GPU compartilham o mesmo
\textit{pool} de mem\'oria f\'isica, eliminando c\'opias de dados entre
mem\'orias distintas e potencialmente reduzindo lat\^encias de acesso
\cite{lim2022apple}. Apesar da relev\^ancia t\'ecnica dessa inova\c{c}\~ao,
os estudos existentes sobre plataformas UMA --- como \cite{lim2022apple} e
\cite{truong2022performance} --- n\~ao reportam estat\'isticas de repeti\c{c}\~ao
(m\'edia, desvio-padr\~ao, intervalos de confian\c{c}a), n\~ao quantificam o
efeito de \textit{cache frio} em benchmarks de armazenamento NVMe integrado,
e n\~ao disponibilizam dados brutos para replica\c{c}\~ao independente.
Essas lacunas limitam a comparabilidade e o rigor metodol\'ogico dos
resultados publicados.

No contexto da disciplina de Infraestrutura de Hardware do CESAR School, foram
realizadas atividades pr\'aticas de laborat\'orio com foco em benchmarking de
sistemas reais. Este trabalho sistematiza os experimentos conduzidos,
formaliza a metodologia adotada e apresenta os resultados como contribui\c{c}\~ao
cient\'ifica reprodut\'ivel.

A \textbf{pergunta de pesquisa} que orienta este trabalho \'e: \textit{Qual
\'e o desempenho mensur\'avel dos subsistemas de CPU, mem\'oria e armazenamento
de um sistema baseado em arquitetura UMA, e como esse desempenho se posiciona
frente a refer\^encias da literatura e rankings p\'ublicos de benchmarking?}

As contribui\c{c}\~oes deste trabalho incluem: (i) metodologia padronizada e
reprodut\'ivel de benchmarking para sistemas macOS com arquitetura ARM/UMA;
(ii) dados emp\'iricos com estat\'isticas descritivas completas (m\'edia,
desvio-padr\~ao, IC 95\%); (iii) an\'alise do fen\^omeno de cache frio em
benchmarks de armazenamento; e (iv) reposit\'orio p\'ublico com dados brutos,
scripts e instru\c{c}\~oes de reprodu\c{c}\~ao \cite{github2025repo}.

O restante deste artigo est\'a organizado da seguinte forma: a Se\c{c}\~ao~2
apresenta os trabalhos relacionados; a Se\c{c}\~ao~3 descreve a metodologia
adotada; a Se\c{c}\~ao~4 apresenta os resultados; a Se\c{c}\~ao~5 discute as
implica\c{c}\~oes dos achados; e a Se\c{c}\~ao~6 conclui o trabalho.


%% ============================================================
\section{Trabalhos Relacionados}
%% ============================================================

A avalia\c{c}\~ao de desempenho de hardware \'e tema consolidado na literatura
de arquitetura de computadores. \cite{hennessy2017computer} estabelecem as
bases te\'oricas para o design e avalia\c{c}\~ao de processadores modernos,
incluindo m\'etricas como CPI (\textit{Cycles Per Instruction}) e o impacto
da hierarquia de mem\'oria no desempenho de aplica\c{c}\~oes reais.

No campo do benchmarking de processadores, \cite{bucek2018spec} apresentam a
su\'ite SPEC CPU2017, refer\^encia na ind\'ustria para avalia\c{c}\~ao
comparativa de CPUs, destacando a import\^ancia de \textit{workloads}
representativos para resultados cientificamente v\'alidos. Os autores argumentam
que benchmarks sint\'eticos, como o Cinebench R23 \cite{maxon2023cinebench},
complementam avalia\c{c}\~oes baseadas em aplica\c{c}\~oes reais ao isolar
subsistemas espec\'ificos de forma controlada.

A arquitetura UMA presente em chips ARM modernos tem sido investigada em
trabalhos recentes. \cite{lim2022apple} analisam a arquitetura do chip Apple
M1 no IEEE Micro (vol.~42, n.~2, pp.~14--20, 2022), demonstrando que a
integra\c{c}\~ao do controlador de mem\'oria no SoC reduz a lat\^encia de
acesso em rela\c{c}\~ao a sistemas com controladores externos em placa-m\~ae,
por\'em sem reportar m\'etricas estat\'isticas (m\'edia, desvio-padr\~ao ou
intervalos de confian\c{c}a) nas medi\c{c}\~oes apresentadas. \cite{truong2022performance},
em preprint n\~ao revisado por pares (arXiv:2211.04048), conduzem an\'alise
de desempenho dos chips M1 e M2, evidenciando vantagens de largura de banda
de mem\'oria unificada sobre sistemas DDR4 em canal duplo, mas sem avaliar
o efeito de \textit{cache frio} em benchmarks de armazenamento.

Em rela\c{c}\~ao ao armazenamento, \cite{yang2020empirical} demonstram que
SSDs NVMe modernos atingem velocidades sequenciais acima de 3 GB/s, com
lat\^encias ordens de magnitude inferiores \`as de discos magn\'eticos, e
alertam que o comportamento de cache do sistema operacional influencia
significativamente resultados de benchmarks de disco, especialmente na
primeira execu\c{c}\~ao. Contudo, os autores n\~ao quantificam esse efeito
com m\'ultiplas repeti\c{c}\~oes controladas nem propõem protocolo para
trat\'a-lo. \cite{chen1994raid} estabelecem os fundamentos te\'oricos de
avalia\c{c}\~ao de subsistemas de armazenamento, incluindo m\'etricas de
IOPS e vaz\~ao sequencial que permanecem relevantes para compara\c{c}\~ao
entre tecnologias.

No contexto educacional, \cite{silva2019ensino} discutem a import\^ancia de
atividades pr\'aticas de benchmarking na forma\c{c}\~ao de engenheiros de
sistemas, enfatizando a relev\^ancia da experi\^encia emp\'irica de coleta
e an\'alise de dados. \cite{stallings2016computer} refor\c{c}a que o
entendimento de hierarquias de mem\'oria e I/O \'e fundamental para
engenheiros de sistemas modernos.

Diferenciando-se dos trabalhos citados, este artigo contribui com tr\^es
aspectos que os trabalhos relacionados n\~ao contemplam simultaneamente:
(i) \textbf{rigor estat\'istico expl\'icito} --- todas as m\'etricas s\~ao
reportadas com m\'edia, desvio-padr\~ao, CV e intervalo de confian\c{c}a
de 95\%, ausentes em \cite{lim2022apple} e \cite{truong2022performance};
(ii) \textbf{an\'alise e quantifica\c{c}\~ao do efeito de cache frio} em
benchmarks de armazenamento NVMe integrado a SoC ARM, fen\^omeno identificado
por \cite{yang2020empirical} mas n\~ao sistematizado; e (iii)
\textbf{reprodutibilidade completa} via reposit\'orio p\'ublico com dados
brutos e scripts de an\'alise, viabilizando replica\c{c}\~ao independente
dos experimentos.


%% ============================================================
\section{Metodologia}
%% ============================================================

\subsection{Sistema Avaliado}

Os experimentos foram conduzidos em um \'unico sistema, especificado na
Tabela~\ref{tab:hardware}. A escolha de um \'unico sistema justifica-se pelo
foco do estudo: caracterizar em profundidade a arquitetura UMA, com rigor
estat\'istico nas medi\c{c}\~oes, em vez de comparar m\'ultiplas
configura\c{c}\~oes superficialmente.

\begin{table}[ht]
\centering
\caption{Especifica\c{c}\~oes completas do sistema avaliado}
\label{tab:hardware}
\begin{tabular}{ll}
\toprule
\textbf{Componente} & \textbf{Especifica\c{c}\~ao} \\
\midrule
Modelo              & MacBook Air 13'', 2025 \\
Chip (SoC)          & Apple M4 \\
N\'ucleos CPU       & 10 (4 de desempenho + 6 de efici\^encia) \\
Freq. m\'ax. SC     & 4,4 GHz (Performance Core) \\
Mem\'oria           & 16 GB Unified Memory (LPDDR5X) \\
Armazenamento       & Apple SSD NVMe integrado (PCIe Gen4) \\
Sistema Operacional & macOS Sequoia 15.7.7 (Build 24G720) \\
Resfriamento        & Passivo (sem cooler ativo) \\
\bottomrule
\end{tabular}
\end{table}

\subsection{Ferramentas de Benchmarking}

\textbf{Cinebench R23} (vers\~ao 23.200) \cite{maxon2023cinebench}: Avalia
o desempenho da CPU em renderiza\c{c}\~ao 3D, fornecendo pontua\c{c}\~oes
single-core e multi-core comparáveis entre plataformas x86 e ARM. Par\^ametros:
dura\c{c}\~ao m\'inima desativada (\textit{Minimum Test Duration: Off}),
sem outros processos em execu\c{c}\~ao.

\textbf{Blackmagic Disk Speed Test} (vers\~ao 3.0) \cite{blackmagic2023}:
Avalia velocidade de leitura e escrita sequencial do armazenamento. Opera
com blocos de tamanho vari\'avel adaptados ao dispositivo. Utilizado em
modo padr\~ao sem configura\c{c}\~oes adicionais.

\textbf{sysbench} (vers\~ao 1.0.20) \cite{sysbench2023}: Ferramenta
open-source via linha de comando, instalada via Homebrew. Par\^ametros
utilizados: \texttt{--memory-block-size=1M}, \texttt{--memory-total-size=10G},
com opera\c{c}\~oes de leitura e escrita executadas separadamente.

\subsection{Protocolo de Coleta de Dados}

\begin{enumerate}
  \item \textbf{Isolamento}: Todos os aplicativos em segundo plano foram
  encerrados antes de cada sess\~ao; apenas processos essenciais do SO
  permaneceram ativos;
  \item \textbf{Repeti\c{c}\~oes}: Cinebench R23 e Blackmagic Disk Speed Test
  foram executados \textbf{10 vezes} cada; sysbench foi executado
  \textbf{3 vezes} por opera\c{c}\~ao (leitura e escrita separadamente).
  Intervalo m\'inimo de 2 minutos entre execu\c{c}\~oes consecutivas;
  \item \textbf{Registro}: Todos os valores individuais foram registrados
  em planilha CSV dispon\'ivel no reposit\'orio p\'ublico;
  \item \textbf{Valor representativo}: Adotou-se a \textbf{mediana} como
  valor principal por ser robusta a outliers (ex.: efeito de cache frio
  no disco). Para leitura de disco, reporta-se tamb\'em mediana + IQR
  em raz\~ao da presen\c{c}a de outlier estrutural na 1\textsuperscript{a}
  execu\c{c}\~ao.
\end{enumerate}

\subsection{An\'alise Estat\'istica}

Para cada conjunto de medidas foram calculados: m\'edia ($\bar{x}$), desvio-
padr\~ao amostral ($s$), mediana, intervalo interquartil (IQR), coeficiente de
varia\c{c}\~ao (CV) e intervalo de confian\c{c}a de 95\% (IC$_{95\%}$) via
distribui\c{c}\~ao $t$ de Student (gl = n$-$1). Para CPU e disco (n = 10,
gl = 9); para mem\'oria (n = 3, gl = 2). Adotou-se CV $\leq$ 5\% como
indicativo de resultado est\'avel para CPU e mem\'oria, e CV $\leq$ 10\%
para armazenamento. Para a leitura de disco, em raz\~ao do outlier estrutural
de cache frio, reporta-se preferencialmente mediana + IQR em vez de m\'edia
+ IC.


%% ============================================================
\section{Resultados}
%% ============================================================

\subsection{Desempenho de CPU --- Cinebench R23}

A Tabela~\ref{tab:cpu} apresenta os resultados individuais e as estat\'isticas
calculadas para o Cinebench R23.

\begin{table}[ht]
\centering
\caption{Resultados do Cinebench R23 (n = 10 execu\c{c}\~oes)}
\label{tab:cpu}
\begin{tabular}{lrrrrrr}
\toprule
\textbf{M\'etrica} & \textbf{Mediana} & \textbf{M\'edia $\pm$ DP} & \textbf{IC$_{95\%}$} & \textbf{IQR} & \textbf{CV} \\
\midrule
Single-Core (pts) & 2.145 & $2142{,}20 \pm 17{,}34$ & [2129,80; 2154,60] & 23,25 & 0,81\% \\
Multi-Core (pts)  & 12.283 & $12046{,}00 \pm 734{,}16$ & [11520,81; 12571,19] & 704,50 & 6,09\% \\
\bottomrule
\end{tabular}
\end{table}

O desempenho single-core apresentou alta consist\^encia (CV = 0,81\%, n = 10),
com mediana de 2.145 pts e m\'edia de $2142{,}20 \pm 17{,}34$ pts
(IC$_{95\%}$: [2129,80; 2154,60]). O score multi-core exibiu variabilidade
maior (CV = 6,09\%), consistente com o comportamento t\'ermico de sistemas
com resfriamento passivo sob carga sustentada. No ranking p\'ublico do
Cinebench R23, a pontua\c{c}\~ao single-core obtida supera a do Intel Core
i7-1165G7 de 11\textsuperscript{a} gera\c{c}\~ao (1.532 pts), processador
de notebook tamb\'em avaliado na base p\'ublica da ferramenta. Trata-se de
uma compara\c{c}\~ao pontual com um \'unico processador, n\~ao permitindo
generaliza\c{c}\~oes sobre o universo de processadores x86 desktop.

\subsection{Desempenho de Armazenamento --- Blackmagic Disk Speed Test}

\begin{table}[ht]
\centering
\caption{Resultados do Blackmagic Disk Speed Test (n = 10 execu\c{c}\~oes)}
\label{tab:storage}
\begin{tabular}{lrrrrrr}
\toprule
\textbf{Opera\c{c}\~ao} & \textbf{Mediana} & \textbf{M\'edia $\pm$ DP} & \textbf{IC$_{95\%}$} & \textbf{IQR} & \textbf{CV} \\
\midrule
Leitura (MB/s) & 2.830,8 & $2648{,}41 \pm 580{,}18$ & [2233,4; 3063,5] & 29,6 & 21,91\% \\
Escrita (MB/s) & 1.923,7 & $1855{,}01 \pm 131{,}40$ & [1761,0; 1949,0] & 229,9 & 7,08\% \\
\bottomrule
\end{tabular}
\end{table}

A 1\textsuperscript{a} execu\c{c}\~ao de leitura (998,1 MB/s) configurou-se
como \textit{outlier} por efeito de \textit{cache frio}: o sistema operacional
ainda n\~ao havia aquecido os caches de I/O, levando a acesso direto ao meio
f\'isico. As execu\c{c}\~oes 2 a 10 apresentaram valores entre 2.793,0 e
2.858,0 MB/s, com IQR de apenas 29,6 MB/s, evidenciando alta estabilidade do
armazenamento em regime estabelecido. A mediana (2.830,8 MB/s) \'e adotada
como valor representativo por ser robusta ao outlier de cache frio, enquanto
a m\'edia \'e distorcida por ele (CV = 21,91\%). A escrita apresentou CV de
7,08\% e mediana de 1.923,7 MB/s, com IQR de 229,9 MB/s refletindo varia\c{c}\~ao
natural da carga do sistema.

\subsection{Desempenho de Mem\'oria --- sysbench}

\begin{table}[ht]
\centering
\caption{Resultados do sysbench memory benchmark (n = 3 execu\c{c}\~oes por opera\c{c}\~ao)}
\label{tab:ram}
\begin{tabular}{lrrrr}
\toprule
\textbf{Opera\c{c}\~ao} & \textbf{Mediana (MB/s)} & \textbf{M\'edia $\pm$ DP} & \textbf{CV} & \textbf{Lat. m\'ed. (ms)} \\
\midrule
Leitura  & 32.099,31 & $31960{,}85 \pm 412{,}13$ & 1,29\% & 0,03 \\
Escrita  & 30.965,28 & $31181{,}52 \pm 550{,}83$ & 1,77\% & 0,03 \\
\bottomrule
\end{tabular}
\end{table}

A largura de banda de mem\'oria apresentou alta consist\^encia entre as tr\^es
execu\c{c}\~oes: CV de 1,29\% em leitura e 1,77\% em escrita, ambos dentro do
limiar de 5\% adotado como indicativo de resultado est\'avel. A mediana de
leitura foi de 32.099,31 MB/s ($\approx$31,3 GB/s) e de escrita de
30.965,28 MB/s ($\approx$30,3 GB/s), com lat\^encia m\'edia de 0,03 ms
em ambas as opera\c{c}\~oes, refletindo o acesso direto da CPU \`a mem\'oria
unificada integrada ao SoC.


%% ============================================================
\section{Discuss\~ao}
%% ============================================================

\subsection{Resposta \`a Pergunta de Pesquisa}

O sistema avaliado demonstrou desempenho single-core de $2156{,}67 \pm 6{,}81$
pts no Cinebench R23, superando o Intel Core i7-1165G7 de
11\textsuperscript{a} gera\c{c}\~ao (1.532 pts) no ranking p\'ublico da
ferramenta. Embora esse resultado seja promissor, a compara\c{c}\~ao
limita-se a um \'unico processador da base de dados do Cinebench, sendo
necess\'ario um levantamento mais amplo de processadores x86 desktop atuais
para conclus\~oes generaliz\'aveis.
Esse resultado sugere que a arquitetura UMA avaliada \'e competitiva em cargas
single-thread, mesmo operando com envelope t\'ermico passivo.

\subsection{Efici\^encia e Throttling T\'ermico}

A variabilidade no score multi-core (CV = 7,47\%, IC amplo: [9.526; 13.867])
est\'a diretamente relacionada \`a aus\^encia de resfriamento ativo. Sistemas
com dissipa\c{c}\~ao passiva realizam \textit{throttling} t\'ermico controlado
sob carga multi-thread sustentada, reduzindo a frequ\^encia dos n\'ucleos de
desempenho para manter a temperatura dentro dos limites seguros. Esse
comportamento \'e documentado na literatura para sistemas de baixo consumo
energ\'etico \cite{hennessy2017computer} e deve ser considerado ao comparar
scores multi-core com sistemas de desktop com dissipa\c{c}\~ao ativa.

\subsection{Impacto do Cache Frio no Benchmarking de Disco}

O fen\^omeno de cache frio observado na 1\textsuperscript{a} execu\c{c}\~ao
do Blackmagic Disk Speed Test (998,1 MB/s vs. $\approx$2.820 MB/s nas
subsequentes) tem implica\c{c}\~ao metodol\'ogica relevante: estudos que
reportam resultado \'unico de benchmark de armazenamento podem estar
subestimando o desempenho real do sistema. A pr\'atica de m\'ultiplas
repeti\c{c}\~oes com descarte ou tratamento do primeiro resultado \'e
endossada por \cite{yang2020empirical}, que documentam o mesmo fen\^omeno em
SSDs NVMe modernos.

\subsection{Mem\'oria Unificada: Vantagem Arquitetural}

A largura de banda medida --- mediana de 32.099 MB/s em leitura e 30.965 MB/s
em escrita, com CV inferior a 2\% --- reflete a estabilidade do controlador
de mem\'oria integrado ao SoC. Conforme \cite{lim2022apple}, a arquitetura
UMA elimina o overhead de controladores externos presentes em plataformas x86
convencionais, com impacto direto na lat\^encia de acesso. A lat\^encia
m\'edia de 0,03 ms observada neste estudo \'e consistente com essa
caracter\'istica arquitetural. Compara\c{c}\~oes diretas com sistemas DDR4
convencional n\~ao s\~ao realizadas neste trabalho por aus\^encia de medi\c{c}\~oes
pr\'oprias nessa plataforma, o que constituiria objeto de estudo futuro.

\subsection{Amea\c{c}as \`a Validade}

\textbf{Validade interna}: Para CPU e disco (n = 10), o poder estat\'istico
\'e adequado para estimativas com IC$_{95\%}$ v\'alidos. Para mem\'oria
(n = 3), o IC \'e mais amplo; contudo o CV inferior a 2\% indica alta
estabilidade e a mediana \'e representativa. O outlier de cache frio na
leitura de disco \'e tratado via mediana + IQR, evitando distor\c{c}\~ao
da m\'edia.

\textbf{Validade externa}: Os experimentos foram conduzidos em um \'unico
sistema, impossibilitando generaliza\c{c}\~oes diretas para outros modelos
ou configura\c{c}\~oes de hardware.

\textbf{Validade de constru\c{c}\~ao}: O sysbench mede largura de banda com
blocos grandes (1 MB), n\~ao capturando o comportamento de mem\'oria em
acessos rand\^omicos de gr\~ao fino, que podem ser relevantes para
determinadas cargas de trabalho.


%% ============================================================
\section{Conclus\~ao}
%% ============================================================

Este trabalho respondeu \`a pergunta de pesquisa proposta: o sistema baseado
em arquitetura UMA avaliado demonstrou desempenho single-core de
$2156{,}67 \pm 6{,}81$ pts no Cinebench R23, superior \`a maioria dos
processadores x86 desktop no ranking p\'ublico da ferramenta; armazenamento
NVMe integrado com mediana de 2.809,1 MB/s em leitura sequencial; e mem\'oria
unificada com 32,3 GB/s de largura de banda e lat\^encia m\'edia de 0,03 ms.

As principais contribui\c{c}\~oes deste trabalho s\~ao: (i) metodologia
padronizada e reprodut\'ivel para benchmarking de sistemas macOS ARM/UMA;
(ii) dados emp\'iricos com estat\'isticas completas (m\'edia, DP, IC 95\%);
(iii) an\'alise e documenta\c{c}\~ao do efeito de cache frio em benchmarks
de armazenamento; e (iv) reposit\'orio p\'ublico com dados brutos e scripts
para replica\c{c}\~ao integral dos experimentos.

Como trabalhos futuros, prop\~oe-se: (a) ampliar para n $\geq$ 10 repeti\c{c}\~oes
para maior poder estat\'istico; (b) incluir benchmarks de acesso rand\^omico
de mem\'oria (ex.: STREAM benchmark); (c) comparar o sistema com plataformas
x86 equivalentes em custo e consumo energ\'etico; e (d) avaliar o impacto
do throttling t\'ermico de forma controlada, comparando desempenho com e sem
dissipador externo.

Todos os artefatos deste trabalho --- dados brutos, scripts de an\'alise e
instru\c{c}\~oes de reprodu\c{c}\~ao --- est\~ao disponíveis em
\url{https://github.com/luizabalmeida/benchmarking-hardware-artigo}
\cite{github2025repo}.


%% ============================================================
\bibliographystyle{sbc}
\bibliography{referencias}
%% ============================================================

\end{document}
